% Applied Software Development Lab -- Lab Handout 8: Things get vaporized
% Copyright (c) Michael Lehn. Licensed under CC BY-SA 4.0.
%
% Region coordinates (bp from the top of the page) follow the page layout of
% the original German handout "Laborblatt 8".
\documentclass{aslab}
\usepackage{unicode-math}
\setmathfont{texgyretermes-math.otf}[Scale=1.2]
\newfontfamily\figserif{TeX Gyre Termes}
\sethandout{8}{Things get vaporized}
\graphicspath{{../figures/handout08/}}

%% inline code in bold lines (the original keeps it in regular weight)
\newcommand*\codemd[1]{{\mdseries\codearg{#1}}}
%% Python listings with the color scheme of the original
\makeatletter
\newcommand*\regionpitch{\baselineskip=\aslab@pitch\relax}
\makeatother
\lstdefinestyle{py}{style=aslab,sensitive=true,escapechar=|,
  basicstyle=\ttfamily\regionpitch,
  morekeywords=[1]{class,def,if,elif,else,for,while,or,and,in,return,not,
    continue},
  keywordstyle=[1]\color{codeorange},
  morekeywords=[2]{import,from},keywordstyle=[2]\color{codemagenta},
  morekeywords=[3]{None,True,False,range,len,int},
  keywordstyle=[3]\color{codecyan},
  morecomment=[l]\#,commentstyle=\color{codepurple},
  alsoletter={0123456789},emphstyle=\color{codered},
  emph={0,1,2,3,4,5,6,7,8,9,01,10,20,30,80,100,105,165,200,223,255,300,400,
    570,600,800}}
\newcommand*\cy[1]{\textcolor{codecyan}{#1}}
%% numbered listing (11bp)
\lstdefinestyle{pyn}{style=py,numbers=left,
  numberstyle=\ttfamily\color{codeorange},numbersep=6.6bp,xleftmargin=26.4bp}
%% numbered listing (10bp)
\lstdefinestyle{pynsmall}{style=py,numbers=left,
  basicstyle=\fontsize{10bp}{11bp}\selectfont\ttfamily\regionpitch,
  numberstyle=\fontsize{10bp}{11bp}\selectfont\ttfamily\color{codeorange},
  numbersep=6bp,xleftmargin=24bp}
%% plain (uncolored) code
\lstdefinestyle{plain}{style=aslab,escapechar=|,
  basicstyle=\ttfamily\regionpitch}

%% lists of this handout
\newlist{ntasks}{enumerate}{2}
\setlist[ntasks]{label={\arabic*.},leftmargin=18bp,labelsep=*,align=left,
  labelindent=0pt}
\newlist{alist}{enumerate}{1}
\setlist[alist]{label={(\Alph*)},leftmargin=18bp,labelsep=*,align=left,
  labelindent=0pt}
\newcommand*\dsh{{\fontsize{13.2bp}{12bp}\selectfont -}}
\newlist{dlist}{itemize}{3}
\setlist[dlist]{label={\dsh},leftmargin=12bp,labelsep=*,align=left,
  labelindent=0pt}
\lineskiplimit=-\maxdimen
\normallineskiplimit=-\maxdimen

\begin{document}
%% ===================================================================== 1 ==
\maketitleblock

\begin{region}{124.0}
\heading{Learning objectives}
\end{region}

\begin{region}[pitch=13]{138.5}
{\color{aslabblue}\bfseries
\begin{dashlist}
\item Deepening object-oriented programming
\item Refactoring existing programs
\item Incremental software development: small changes, testing, and versioning
\end{dashlist}\par}
\end{region}

\begin{region}{191.5}
\heading{Step 1: Cloning your own repository}
\end{region}

\begin{region}[pitch=13.17]{219.1}
Today, too, we are working ``only'' with Python. The goal is to add new
features. For example, balls that you miss with the paddle should disappear at
the bottom edge. And -- as you would expect from Breakout -- there should be
stones (bricks) that disappear when the ball hits them. The ball is reflected
in the process.
\end{region}

\begin{region}[pitch=13.3]{285.2}
Our project currently consists of only a few modules (\code{ball.py},
\code{breakout.py}, \code{controller.py}, \code{game.py}, \code{paddle.py},
\code{rocket.py}, \code{serial-reader.py}). Still, you should not underestimate how
hard it is to keep track of a program that consists of several parts -- even
if each file is only a few lines long.
\end{region}

\begin{region}{351.5}
You need
\end{region}

\begin{region}[pitch=14.3]{364.6}
\begin{dashlist}
\item on the one hand, an eye for the \textbf{big picture}: What does each
  module do?
\item on the other hand, an eye for \textbf{detail}: How does the module do
  what it is supposed to do?
\end{dashlist}
\end{region}

\begin{region}[pitch=13.15]{406.4}
You have to switch flexibly back and forth between these two levels. Beyond a
certain level of complexity, keeping both in mind at the same time is neither
possible nor necessary. This is the skill we want to train in the first step.
\end{region}

\begin{region}[pitch=13.2]{459.0}
We proceed as follows: start the game, observe what it can do, then skim the
code to recapitulate its structure.
\end{region}

\begin{region}[pitch=13.3]{499.5}
\subheading{Preparation: Opening two terminals}
As always: left terminal for editing, right terminal for running commands.
\end{region}

\begin{region}[pitch=14.3]{540.1}
\subheading{Right terminal (cloning the repository and starting the game)}
First, clone your repository with
\end{region}

\begin{region}{567.9}
\code{git clone https://gitlab.uni-ulm.de/USERNAME/asp-ulm}
\end{region}

\begin{region}[pitch=13.5]{581.1}
Use \code{cd asp-ulm/python} to change into the Python directory, and then
start the game with \code{python breakout.py}.
\end{region}

\begin{region}{621.1}
You should observe:
\end{region}

\begin{region}[pitch=14]{633.7}
\begin{dashlist}
\item The balls are cleanly reflected at all edges (top, bottom, left, right).
  More precisely: the \textbf{center} of the balls always stays within the
  drawing surface -- this was taken into account in the reflection.
\end{dashlist}
\end{region}

\begin{region}[pitch=15.3]{674.8}
\begin{dashlist}
\item You can make the ball bounce off the \textbf{paddle}. But ``angle of
  incidence = angle of reflection'' does \textbf{not} apply here:
\end{dashlist}
\end{region}

\begin{region}[x=68.7,right=0,pitch=14.25]{702.8}
\begin{dashlist}
\item If you hit the ball in the middle, it flies straight up.
\item Far left $\rightarrow$ 30° to the left.
\item Far right $\rightarrow$ 30° to the right.
\end{dashlist}
\end{region}

\begin{region}[pitch=14.3]{745.8}
\begin{dashlist}
\item With the \textbf{up arrow key}, the paddle fires rockets (small discs)
  that fly straight up and disappear at the top edge.
\end{dashlist}
\end{region}

\nextpage
%% ===================================================================== 2 ==
\begin{region}[pitch=14.3]{69.4}
\subheading{Left terminal (examining the code)}
Now the point is to understand \textbf{how} these features were implemented.
There is no ``right'' or ``wrong'' here -- that depends on style, taste, and
also on the state of the project.
\end{region}

\begin{region}{123.4}
One possible way is the top-down approach:
\end{region}

\begin{region}[pitch=13.37]{136.8}
\begin{tasks}
\item Start with the main loop in \code{breakout.py}.
\item Among other things, you will see: there are classes such as \code{Ball}
  and \code{Rocket}, of which several instances are created, each with a
  \code{draw} and a \code{move} method.
\item There is also the class \code{Paddle}, of which there is only one
  instance -- with (static) methods such as \code{left}, \code{right}, and
  \code{draw}.
\item Then you can go through the files \code{ball.py}, \code{rocket.py}, and
  \code{paddle.py} one by one.
\item While reading, keep asking yourself questions such as:
  \begin{bullets}
  \item ``Why did I do it this way back then?''
  \item ``What exactly does this or that mean in Python again?''
  \item ``How are the modules connected to each other?''
  \end{bullets}
\end{tasks}
\end{region}

\begin{region}[pitch=13.93]{283.7}
You may notice: some modules such as \code{controller.py} or
\code{serial_reader.py} are \textbf{not used at all} when the game starts.
This is not a mistake -- these files concern the \textbf{serial connection},
which will later allow you to control the game with a \textbf{game
controller} instead of the keyboard.
\end{region}

\begin{region}[pitch=14.25]{352.8}
Today, however, it is \textbf{only} about the game features. It is perfectly
sufficient if you roughly recall \emph{what} these modules are meant for. You
do not need to look at them in detail today.\\
\textbf{One thing at a time.}
\end{region}

\begin{region}[pitch=13.23]{409.8}
\subheading{Task}
Go through this process on your own. Start the game, observe its behavior, and
examine the code.\\
If anything is unclear -- ask!
\end{region}

\begin{region}[pitch=13.3]{490.0}
\heading{Remark}
While looking through the code, we came across a variant like this:
\end{region}

\begin{region}[pitch=13.2]{529.9}
\begin{lstlisting}[style=plain]
for r in rockets:
    r.draw()
    r.move()
    rockets = [item for item in rockets if item.alive]
\end{lstlisting}
\end{region}

\begin{region}[pitch=14.28]{596.6}
In Python, this \textbf{seemingly works without problems}, even though the
list \code{rockets} is \textbf{modified during the iteration}. The reason: in
each pass of the loop, a \textbf{new list} is created and the name
\code{rockets} is rebound -- but the loop \textbf{keeps running over the
original iterator}, which still refers to the old list. So it only works
because Python allows \textbf{inefficient (and potentially confusing)} code.
\end{region}

\begin{region}[pitch=13.2]{680.4}
\begin{lstlisting}[style=plain]
for r in rockets:
    r.draw()
    r.move()
rockets = [item for item in rockets if item.alive]
\end{lstlisting}
\end{region}

\begin{region}[pitch=13.6]{747.1}
Here, \textbf{all rockets are first moved and drawn}, and \textbf{then} all
rockets that are no longer visible (\code{alive == False}) are removed.
\end{region}

\begin{region}{778.3}
This is \textbf{more efficient} and \textbf{easier to understand} -- also for
your future self in a few weeks \emoji{1F609}
\end{region}

\nextpage
%% ===================================================================== 3 ==
\begin{region}{69.4}
\heading{Step 2: Balls you miss \dots}
\end{region}

\begin{region}{97.0}
Now we test whether recapitulating in the first step has paid off.
\end{region}

\begin{region}[pitch=13.4]{124.3}
\textbf{Goal:} If you miss a ball with the paddle, it should disappear as soon
as it reaches the bottom edge of the drawing surface. Among other things, this
means: the corresponding ball object should be removed from the list
\code{balls}.
\end{region}

\begin{region}[pitch=13.4]{178.4}
\subheading{Task}
Implement this feature:
\end{region}

\begin{region}[pitch=13.4]{204.4}
\begin{dashlist}
\item Detect when a ball is ``lost'' (e.g., based on its y-position).
\item Remove the corresponding object from the list \code{balls}.
\end{dashlist}
\end{region}

\begin{region}[pitch=14.25]{232.6}
\textbf{Tip:} Think about \textbf{how this was done for the rockets}. There,
too, objects disappear when they reach the edge. The goal is to understand
this mechanism and \textbf{transfer it to the balls in the same way}.
\end{region}

\begin{region}[pitch=14.3]{288.5}
If you have \textbf{questions}, let us know. If you think everything works,
\textbf{show us} what you have changed. \textbf{Once we give you the OK},
off to the repository:
\end{region}

\begin{region}[pitch=13.2]{316.4}
\code{git commit -m "ball disappears"}\\
\code{git push}
\end{region}

\nextpage
%% ===================================================================== 4 ==
\begin{region}{69.4}
\heading{Step 3: Space Invaders in Breakout?}
\end{region}

\begin{region}[pitch=13.2]{97.0}
Why not? This is an invasion in the truest sense of the word -- and at the
same time a good moment to notice how much it pays off to have looked at the
existing code beforehand.
\end{region}

\begin{region}[pitch=13.25]{150.7}
\textbf{Idea:} Analogous to the class \code{Ball}, you can now program a
class \code{Ship}.\\
Each object of this class represents a spaceship that moves across the
drawing surface -- just like the ``invaders'' in the classic game.
\end{region}

\begin{region}[pitch=14.0]{204.4}
\subheading{Movement along a curve}
Of course, the spaceships should \textbf{not simply move in a straight line}
-- that would make them too easy to hit. Instead, you can let them fly along
an interesting curve. A nice option is the so-called \textbf{lemniscate of
Gerono}. The trace of this curve looks like a \textbf{figure eight lying on
its side}:
\end{region}

\placefigure{56.7}{264.4}{477.4}{lemniscate.tex}

\begin{region}{642.8}
The curve is parametrized by time (with $t$ as the time) and reads:
\end{region}

\begin{region}{673.1}
\centering $x(t) = x_0 + a\cdot\cos t,\quad y(t) = y_0 + b\cdot\sin 2t$
\end{region}

\begin{region}{700.9}
This means: the spaceship moves within a rectangle with the corners:
\end{region}

\begin{region}[pitch=18.37]{717.1}
\begin{dashlist}
\item \textbf{Top left:} $(x_0-a, y_0-b)$
\item \textbf{Bottom left:} $(x_0-a, y_0+b)$
\item \textbf{Bottom right:} $(x_0+a, y_0+b)$
\item \textbf{Top right:} $(x_0+a, y_0-b)$
\end{dashlist}
\end{region}

\nextpage
%% ===================================================================== 5 ==
\begin{region}{69.4}
\subheading{Task: Bringing spaceships into the game}
\end{region}

\begin{region}[pitch=13.3]{89.6}
Write a new class \code{Ship} -- analogous to the class \code{Ball}. A
spaceship moves along a lemniscate and can (later) be shot down.
\end{region}

\begin{region}[pitch=13.5]{123.6}
\subheading{Part 1: Creating the class Ship}
Create a new file \code{ship.py} and define a class \code{Ship} in it.\\
Since this class is very similar to the class \code{Ball}, you can simply
start with a copy:
\end{region}

\begin{region}{170.7}
\code{cp ball.py ship.py}
\end{region}

\begin{region}{190.6}
Then modify the file step by step:
\end{region}

\begin{region}{211.3}
\begin{alist}[label={\bfseries(\Alph*)}]
\item \textbf{Constructor}
\end{alist}
\end{region}

\begin{region}[x=74.7,right=0,pitch=13.4]{224.6}
A spaceship should be created like this:
\begin{lstlisting}[style=plain]
ship = Ship(0)
\end{lstlisting}
\end{region}

\begin{region}[x=74.7,right=0]{254.6}
The parameter \code{0} stands for the initial value of the time $t=0$.
\end{region}

\begin{region}[x=74.7,right=0,pitch=13.45]{275.0}
The class should have the following attributes:
\begin{lstlisting}[style=plain]
self.x0 = Game.width / 2
self.y0 = Game.height / 3
self.a = 0.9 * Game.width / 2
self.b = 0.9 * Game.height / 3
\end{lstlisting}
\end{region}

\begin{region}[x=74.7,right=0,pitch=13.35]{355.4}
In addition, you need:
\begin{lstlisting}[style=plain]
self.t = t                  # current value of the time
self.alive = True           # False means shot down
self.x, self.y = None, None # current position (set in move)
self.move()                 # call move to set the position
\end{lstlisting}
\end{region}

\begin{region}{429.2}
\begin{alist}[start=2,label={\bfseries(\Alph*)}]
\item \textbf{Method \codemd{move(self)}}
\end{alist}
\end{region}

\begin{region}[x=74.7,right=0,pitch=13.25]{442.5}
Here you compute the new position along the lemniscate:
\begin{lstlisting}[style=py]
def |\cy{move}|(self):
    self.x = self.x0 + self.a * math.cos(self.t)
    self.y = self.y0 + self.b * math.sin(2 * self.t) / 2
    self.t += 0.01
    while self.t > 2 * math.pi:
        self.t -= 2 * math.pi
\end{lstlisting}
\end{region}

\begin{region}[x=74.7,right=0,pitch=17.2]{538.2}
Resetting $t$ to the interval $[0,2\pi)$ is \textbf{not strictly necessary} --
Python has no problem with numbers growing arbitrarily large (and don't worry:
in this case, the memory will not fill up in your lifetime \emoji{1F604}).
But: it is more elegant and saves memory -- and it makes a true engineer's
heart beat faster \emoji{2705}
\end{region}

\begin{region}{613.0}
\begin{alist}[start=3,label={\bfseries(\Alph*)}]
\item \textbf{Method \codemd{draw(self)}}
\end{alist}
\end{region}

\begin{region}[x=74.7,right=0,pitch=13.2]{626.3}
Simply draw the spaceship as a disc with a larger radius than the ball -- and
preferably in a different color.
\end{region}

\begin{region}[pitch=13.6]{666.7}
\subheading{Part 2: Building the spaceships into the game}
Now the spaceships have to be integrated into \code{breakout.py} -- analogous
to the balls:
\end{region}

\begin{region}[pitch=13.4]{693.8}
\begin{ntasks}
\item Importing the class:\ \ \code{from ship import Ship}
\item Creating a list of ships, e.g., \code{ships = [Ship(0), Ship(0.5), Ship(1)]}
\item Moving, drawing, and removing ships that have been shot down from the
  list
\end{ntasks}
\end{region}

\begin{region}[pitch=13.5]{740.3}
Once your code runs:
\end{region}

\begin{region}[pitch=13.2]{753.8}
\code{git add ship.py}\\
\code{git commit -a -m "spaceships!"}\\
\code{git push}
\end{region}

\nextpage
%% ===================================================================== 6 ==
\begin{region}{69.4}
\heading{Step 4: Now it's getting serious!}
\end{region}

\begin{region}[pitch=13.7]{90.3}
Of course, you should be able to \textbf{shoot down a spaceship} with a
\textbf{ball} or a \textbf{rocket}. This works by checking whether the
distance between a spaceship and a ball or a rocket is \textbf{small enough}.
\end{region}

\begin{region}{138.6}
Here is an example of how to check whether a \textbf{ball} has come too close
to a \textbf{spaceship}:
\end{region}

\begin{region}{152.2}
\begin{lstlisting}[style=py]
    for s in ships:
\end{lstlisting}
\end{region}

\begin{region}{165.4}
\begin{lstlisting}[style=py]
        for b in balls:
            dist = math.hypot(b.x - s.x, b.y - s.y)
            if dist <= 20:
                s.alive = False
\end{lstlisting}
\end{region}

\begin{region}[pitch=14.3]{220.6}
\subheading{What exactly happens here?}
We want to find out whether a \textbf{ball has hit a spaceship}. To do so, we
have to calculate, for \textbf{every spaceship}, the distance to \textbf{every
ball}. For this purpose, two loops are \textbf{nested inside each other}:
\end{region}

\begin{region}[pitch=14.4]{276.8}
\begin{dashlist}
\item The \textbf{outer loop} runs with the variable \code{s} over all
  spaceships in the list \code{ships}.
\item For \textbf{each individual value of} \code{s}, the \textbf{inner loop}
  runs with the variable \code{b} over all balls in the list \code{balls}.
\end{dashlist}
\end{region}

\begin{region}[pitch=13.23]{319.8}
So in the \textbf{innermost loop}, there is always exactly \textbf{one
spaceship} \code{s} and one ball \code{b} -- for every possible pair of
spaceship and ball.\\
Then the distance between these two points is calculated with good old
Pythagoras:
\end{region}

\begin{region}{381.2}
\centering $\text{dist} = \sqrt{(\text{b.x} - \text{s.x})^2
  + (\text{b.y} - \text{s.y})^2}$
\end{region}

\begin{region}[pitch=13.2]{407.1}
If the distance is less than or equal to 20, the spaceship has been hit. Why
20? In my version, a ball has radius 5 and a spaceship has radius 15.
\end{region}

\begin{region}{447.6}
\subheading{Task:}
\end{region}

\begin{region}{468.4}
\begin{ntasks}
\item \textbf{Import} the module \code{math} in \code{breakout.py}, if you
  have not done so yet:
\end{ntasks}
\end{region}

\begin{region}[x=74.7,right=0]{488.6}
\begin{lstlisting}[style=py]
import math
\end{lstlisting}
\end{region}

\begin{region}[x=74.7,right=0]{509.4}
\textbf{Insert the above code} at a suitable place in the main loop:
\end{region}

\begin{region}[x=74.7,right=0,pitch=14.25]{523.7}
\begin{bullets}
\item \textbf{after} all objects have been moved and drawn
\item \textbf{before} the lists \code{balls} and \code{ships} are cleaned up
  (i.e., the ``corpses'' are removed)
\end{bullets}
\end{region}

\begin{region}[x=74.7,right=0]{552.2}
\textbf{Test} whether you can now shoot down a spaceship with a ball.
\end{region}

\begin{region}[pitch=13.85]{565.5}
\begin{ntasks}[start=2]
\item Actually, the point of the rockets was to be able to shoot something
  down with them.\\
  Change the code so that the inner loop iterates over the list
  \code{rockets} instead of \code{balls}. \textbf{Test} whether you can now
  shoot down a spaceship with a rocket.
\end{ntasks}
\end{region}

\begin{region}[pitch=13.25]{607.4}
\begin{ntasks}[start=3]
\item It would of course be even better if you could shoot down a spaceship
  \textbf{with both a ball and a rocket} -- without duplicate code.\\
  An obvious approach would be, for example, the following:
\end{ntasks}
\end{region}

\begin{region}[x=74.7,right=0,pitch=13.2]{653.9}
\begin{lstlisting}[style=py]
      for b in [balls, rockets]:
            |…|
\end{lstlisting}
\end{region}

\begin{region}[x=74.7,right=0,pitch=14.3]{687.6}
You might think that this iterates over \textbf{all balls and all rockets}.
\textbf{But convince yourselves that this does not work!}
\end{region}

\begin{region}[x=74.7,right=0,pitch=13.5]{729.3}
\subheading{So what really happens here?}
\code{[balls, rockets]} is a list that contains exactly two elements: namely
the list \code{balls} and the list \code{rockets}. This means: the loop runs
twice -- once with \code{b = balls}, once with \code{b = rockets} -- and
\code{b} is then a whole list each time, not a single object.
\end{region}

\nextpage
%% ===================================================================== 7 ==
\begin{region}[pitch=13.5]{68.8}
\begin{ntasks}[start=4]
\item With \code{balls + rockets}, you create a new list that contains all
  elements of \code{balls} and of \code{rockets}.
\end{ntasks}
\end{region}

\begin{region}[x=74.7,right=0]{108.7}
So if you write:
\end{region}

\begin{region}[x=74.7,right=0,pitch=13.2]{135.3}
\begin{lstlisting}[style=py]
for s in ships:
    for b in balls + rockets:
        dist = math.hypot(b.x - s.x, b.y - s.y)
        if dist <= 20:
            s.alive = False
\end{lstlisting}
\end{region}

\begin{region}[x=74.7,right=0,pitch=17.5]{214.2}
then everything should work as intended: you can vaporize a spaceship with a
ball as well as with a rocket \emoji{1F4A5}
\end{region}

\begin{region}[x=74.7,right=0,pitch=13.2]{260.6}
Test whether this works. If so -- then you know that the effort has paid off.
And then, as always, into the repository with it:
\end{region}

\begin{region}[x=74.7,right=0,pitch=13.5]{300.4}
\code{git commit -a -m "shoot to kill"}\\
\code{git push}
\end{region}

\begin{region}{341.0}
\heading{Step 5: Bricks!}
\end{region}

\begin{region}[pitch=17.7]{369.5}
Now a \textbf{wall of stones (bricks)} enters the game -- which of course
makes sense if the game is called \textbf{Breakout} \emoji{1F609}
\end{region}

\begin{region}[pitch=14.3]{403.9}
In a \textbf{first version}, the balls should simply still \textbf{fly
through the bricks}. At first, the point is \textbf{not} the collision, but
rather:
\end{region}

\begin{region}{435.8}
\emoji{1F9E9} \textbf{Where on the drawing surface is there a brick -- and
where not?}
\end{region}

\begin{region}[pitch=14.03]{466.8}
\subheading{The grid}
For simplicity, a \textbf{brick} in our case is \textbf{80 pixels wide} and
\textbf{20 pixels high}. Since the drawing surface is \textbf{800 pixels
wide} and \textbf{600 pixels high}, this results in a \textbf{grid with 10
columns and 30 rows}. Of course, not every cell in this grid is occupied by a
brick.\\
For example, only a few of the rows could contain bricks at first -- as in
the following sketch:
\end{region}

\placefigure{56.7}{526.8}{481.9}{brick-grid.tex}

\nextpage
%% ===================================================================== 8 ==
\begin{region}[pitch=14.3]{82.6}
\subheading{A bit of mathematics}
A \textbf{grid} with \textbf{10 columns} and \textbf{30 rows}, in which each
cell either contains a brick or not, can be thought of as a
\textbf{matrix}.\\
Each entry is then either \code{1} $\rightarrow$ brick present, or \code{0}
$\rightarrow$ no brick present. And because we, as mathematicians, value
precision and conventions: we speak of a \textbf{30\,×\,10 matrix} (rows ×
columns).
\end{region}

\begin{region}[pitch=13.6]{175.2}
\textbf{Implementation in Python: \codemd{numpy}}\\
In Python, the module \code{numpy} lets you work with matrices very
efficiently. For our purposes, we only need:
\end{region}

\begin{region}[pitch=14.3]{215.2}
\begin{dashlist}
\item a \textbf{zero-initialized 30\,×\,10 matrix}.
\item \textbf{Indexing} starts \textbf{at 0}: the row index runs from 0 to 29,
  the column index from 0 to 9.
\item access to individual entries for \textbf{reading and writing}.
\end{dashlist}
\end{region}

\begin{region}[pitch=13.6]{265.6}
\textbf{The class \codemd{Wall}}\\
Create a new file \code{wall.py} and write the following code in it:
\end{region}

\begin{region}[pitch=13]{293.1}
\begin{lstlisting}[style=pynsmall]
import pygame
import numpy
from game import Game

class |\cy{Wall}|:
    brick = numpy.zeros((30, 10))

    def |\cy{init}|():
        for i in range(4, 8):
            for j in range (0, 10):
                Wall.brick[i, j] = 1

    def |\cy{draw}|():
        for i in range(0, Wall.brick.shape[0]):
            for j in range(0, Wall.brick.shape[1]):
                if Wall.brick[i, j] == 0:
                    continue
                w = 80
                h = 20
                x = j * w
                y = i * h
                r = (x + 1, y + 1, w - 1, h - 1)
                pygame.draw.rect(Game.screen, (223, 105, 165), r)
\end{lstlisting}
\end{region}

\begin{region}{599.9}
\textbf{Explanation of \codemd{wall.py}}
\end{region}

\begin{region}[pitch=13.65]{613.3}
\begin{dlist}
\item In \textbf{line 2}, the module \code{numpy} is imported.
\item The class has a \textbf{class attribute} \code{brick}, which is a
  \textbf{30\,×\,10 matrix full of zeros}.
\item The method \code{init} sets all matrix entries in \textbf{rows 4 to 7}
  to \code{1} -- that is where the bricks are!
\item The method \code{draw} draws the wall:
  \begin{dlist}
  \item With \code{Wall.brick.shape[0]} and \code{Wall.brick.shape[1]}, you
    query the number of rows and the number of columns, respectively.
  \item Two nested loops run through the matrix.
  \item \code{i} stands for the \textbf{row index}, \code{j} for the
    \textbf{column index}.
  \item If the entry is \code{0}, \code{continue} jumps to the next loop
    iteration.
  \item Otherwise, a rectangle is drawn:
    \begin{dlist}
    \item width \code{w = 80}, height \code{h = 20}
    \item top left corner at \code{x = j * w}, \code{y = i * h}
    \item slightly shifted (\code{+1}) and reduced in size (\code{-2}) for
      visible gaps
    \end{dlist}
  \end{dlist}
\end{dlist}
\end{region}

\nextpage
%% ===================================================================== 9 ==
\begin{region}{69.4}
\subheading{Tasks:}
\end{region}

\begin{region}[pitch=13.9]{96.9}
\begin{ntasks}
\item \textbf{Build the brick wall into the game} by integrating the class
  Wall into \code{breakout.py}:
  \begin{alist}
  \item Import the class at the beginning of the file:
    \code{from wall import Wall}
  \item After \code{Game.init()}, also call the initialization of the wall:
    \code{Wall.init()}
  \item Within the main loop, also call \code{Wall.draw()} -- preferably
    \textbf{after the balls have been drawn}: Wall.draw()
  \end{alist}
\end{ntasks}
\end{region}

\begin{region}[pitch=13.45]{166.7}
\begin{ntasks}[start=2]
\item \textbf{Test} whether the brick wall is displayed as expected.\\
  If everything works -- into the repository:\\
  \code{git add wall.py}\\
  \code{git commit -a -m "the wall is up"}\\
  \code{git push}
\end{ntasks}
\end{region}

\begin{region}{260.7}
\heading{Step 6: Now we clear the bricks out of the way!}
\end{region}

\begin{region}[pitch=14.1]{287.2}
The bricks are currently pure decoration. In this step, they are to fulfill
their purpose: \textbf{they are to be smashed by the balls.}
\end{region}

\begin{region}[pitch=13.25]{328.8}
Of course, one could calculate exactly \textbf{whether and where a ball hits a
brick}: after all, for each frame we know the old and the new position of a
ball. From this, one could calculate which edge of a brick is crossed.
\end{region}

\begin{region}{372.8}
\emoji{1F4A1} Anyone who feels like a little math puzzle is welcome to
``solve for x'' on their own.
\end{region}

\begin{region}{402.6}
We make it easier for ourselves at this point -- with an
\textbf{approximation}:
\end{region}

\begin{region}[pitch=14.3]{422.6}
\begin{dashlist}
\item From the \textbf{old position} \code{x1}, \code{y1} of the ball, we
  calculate \textbf{in which grid cell} the ball was before: this gives a pair
  \code{i1}, \code{j1} (row, column).
\item From the \textbf{new position} \code{x2}, \code{y2}, we calculate
  accordingly which \textbf{grid cell} the ball has now entered: this gives a
  pair \code{i2}, \code{j2} (row, column).
\end{dashlist}
\end{region}

\begin{region}[pitch=14.25]{485.7}
If the ball moves \textbf{from one cell into another cell with a brick}, we
assume that a brick was smashed in the process -- and \textbf{set the entry in
the matrix} \code{Wall.brick[i2, j2]} to 0.
\end{region}

\begin{region}{541.6}
\textbf{And there is more:} We can even estimate approximately \textbf{how}
the brick was hit:
\end{region}

\begin{region}[pitch=13.8]{555.0}
\begin{dashlist}
\item If \code{i1 ≠ i2}, the \textbf{top or bottom edge} was probably smashed
  through.
\item If \code{j1 ≠ j2}, the \textbf{left or right edge} was probably smashed
  through.
\item If both \code{i1 ≠ i2} and \code{j1 ≠ j2} hold, a corner would even
  have to have been hit, which, however, is rather rare in practice. Such cases
  also clearly show the limits of this approximation.
\end{dashlist}
\end{region}

\begin{region}[pitch=13.3]{637.5}
\textbf{The following sketch shows three typical cases} -- and helps you get
a feeling for how well (or how poorly) this method works in practice:
\end{region}

\placegraphic{56.7}{654.6}{481.9}{brick-cases.png}

\nextpage
%% ==================================================================== 10 ==
\begin{region}{69.4}
\textbf{How do you calculate the column and row in the grid from the
(x,\,y) coordinates?}
\end{region}

\begin{region}[pitch=14.6]{89.4}
\begin{dashlist}
\item If you divide the \textbf{x-coordinate} by the \textbf{width} of a cell
  and round the result down to an integer, you get the \textbf{column}. In
  Python, this is easily done with integer division (\code{//}).
\item The same applies to the \textbf{y-coordinate}: if you divide it by the
  \textbf{height} of a cell and round down, you get the \textbf{row}.
\end{dashlist}
\end{region}

\begin{region}{153.0}
Add the following class method to the class Wall:
\end{region}

\begin{region}{173.0}
\begin{lstlisting}[style=pyn,firstnumber=13]
    def |\cy{collision}|(x2, y2):
\end{lstlisting}
\end{region}

\begin{region}{186.2}
\begin{lstlisting}[style=pyn,firstnumber=14]
        j2 = int(x2) // 80
        i2 = int(y2) // 20
        if i2 >= 0 and i2 < 30 and j2 >= 0 and j2 < 10:
            if Wall.brick[i2, j2]:
                Wall.brick[i2, j2] = 0
\end{lstlisting}
\end{region}

\begin{region}{252.9}
\subheading{What does this method do?}
\end{region}

\begin{region}[pitch=13.75]{265.0}
\begin{dashlist}
\item In this version, \code{collision} is initially passed only the
  \textbf{new position} \code{(x2, y2)} of the ball.
\item In \textbf{lines 14 and 15}, the column \code{j2} and the row \code{i2}
  are calculated from \code{x2} and \code{y2} -- by integer division by the
  cell width and cell height, respectively.
\item In \textbf{line 16}, we check whether \code{i2} and \code{j2} are valid
  indices in the grid at all -- because the ball may already be
  \textbf{outside} the drawing surface (e.g., before the correction at the
  walls).
\item In \textbf{line 17}, we check whether \textbf{there is a brick in the
  cell that was hit}.
\item In \textbf{line 18}, this brick is then removed -- so we ``vaporize''
  the brick by setting the matrix entry to \code{0}.
\end{dashlist}
\end{region}

\begin{region}{411.2}
\subheading{Tasks:}
\end{region}

\begin{region}[pitch=13.5]{431.3}
\begin{ntasks}
\item In the \code{move} method of the class \code{Ball}, we now want to check
  whether the ball hits a brick. Only a few changes in \code{ball.py} are
  needed for this:
  \begin{alist}
  \item The class \code{Wall} has to be imported: \code{from wall import Wall}
  \item In the \code{move} method (after \code{new_x} and \code{new_y} have
    been calculated and corrected), call the collision detection:
    \code{Wall.collision(new_x, new_y)}
  \end{alist}
\end{ntasks}
\end{region}

\begin{region}[x=74.7,right=0,pitch=13.2]{499.6}
\textbf{Test whether bricks now disappear on contact with a ball.} The ball
should fly through the bricks like a knife through butter -- because its
course currently remains completely unaffected. We only check whether it hits
a brick -- and then remove that brick.
\end{region}

\begin{region}[pitch=13.35]{540.2}
\begin{ntasks}[start=2]
\item \textbf{Now comes today's grand finale.}\\
  On contact with a brick, the ball should not only make the brick disappear
  but also bounce off. To do this, change the call in the \code{move} method
  of the class \code{Ball} to:
\end{ntasks}
\end{region}

\begin{region}[x=74.7,right=0,pitch=13.2]{580.4}
\begin{lstlisting}[style=plain]
self.vx, self.vy = Wall.collision(self.x, self.y, new_x, new_y,
                                  self.vx, self.vy)
\end{lstlisting}
\end{region}

\begin{region}[x=74.7,right=0,pitch=13.4]{607.1}
Now the old position (\code{self.x}, \code{self.y}), the new position
(\code{new_x}, \code{new_y}), and the current velocity vector
(\code{self.vx}, \code{self.vy}) are passed.\\
The method should now return a (possibly adjusted) velocity vector. Now it is
your turn to extend \code{Wall.collision(...)} accordingly:
\end{region}

\begin{region}[x=74.7,right=0,pitch=13.6]{661.5}
\begin{ntasks}
\item If the \textbf{top or bottom edge} of a brick is smashed through, the
  \textbf{y-component} of the velocity vector should be ``reversed''.
\item If the \textbf{left or right edge} of a brick is smashed through, the
  \textbf{x-component} of the velocity vector should be ``reversed''.
\item If \textbf{both indices} change, the ball bounces off a corner -- and
  \textbf{both components} should be negated.
\end{ntasks}
\end{region}

\begin{region}[x=74.7,right=0]{744.0}
Test whether the ball now bounces off correctly and removes the brick. If it
all works:
\end{region}

\begin{region}[x=74.7,right=0,pitch=13.2]{757.4}
\code{git commit -a -m "bricks can take a hit now"}\\
\code{git push}
\end{region}

\end{document}
